% p3-03-trinity-monomials

\section{P3 §3. Multiplicative Grammar of Trinity Monomials}

Figure (fig-p3-ch3-trinity-grammar). Trinity grammar triptych (Paper III recap) --- basis abacus,
                              open (p,m) lattice, five-spoke generator wheel.

  3.1 The Trinity Logarithmic Group
  The set T under multiplication forms an abelian monoid; passing to logarithms defines

  L = { log T : T $\in$ T } = { log n + k log 3 + p log $\varphi$ + m log pi + q log e : n $\geq$ 1, (k,p,m,q)
  $\in$ Z4 }.

  This is the additive group generated by G = { log 3, log $\varphi$, log pi, 1 } over Z, extended
  by the discrete rational scaling dimension from n. The generators are conjecturally
  Q-linearly independent (Working hypothesis; see Section 11.1).

        Proposition 3.1 (Injectivity of Trinity label map; conditional). If log 3, log $\varphi$,
        log pi, and 1 are linearly independent over Q, then the map lambda ↦ T(lambda)
        from Z$\geq$ 1 $\times$ Z4 to T is injective.

  Proof sketch. Suppose T(lambda1) = T(lambda2). Taking logarithms yields a Q-linear
  relation among log 3, log $\varphi$, log pi, 1; by the independence assumption this forces
  lambda1 = lambda2. □

        Remark. Baker's theorem on linear forms in logarithms of algebraic numbers
        covers log 3 and log $\varphi$ (both algebraic, hence Baker's theorem applies), but log pi
        involves a transcendental argument not covered by the standard machinery. The
        full independence is therefore adopted as a Working hypothesis throughout; see
        Section 11 for a complete discussion.

        Remark (Anchor embedding). The structural observation log 3 = log($\varphi$2 + $\varphi$-2)
        from (star) illuminates why log 3 and log $\varphi$ are algebraically distinct: the logarithm
        of a sum is not the sum of logarithms, so even though both $\varphi$2 and $\varphi$-2 are
        algebraic, log 3 and log $\varphi$ remain linearly independent over Q by Baker's theorem.
        This structural independence supports the working independence hypothesis in

        the algebraic part of the generator set.

  3.2 Partial Order by Complexity
  Define a partial order on T by T1 preceq T2 iff ell(lambda1) $\leq$ ell(lambda2). The
  minimal elements are integers T(n,0,0,0,0) = n, followed by monomials with a single
  nonzero exponent.

        Definition 3.2 (Complexity shell SL). For integer L $\geq$ 0,

  SL = { T $\in$ T : ell(lambdaT) = L }.

  The growth rate |SL| $\sim$ A $\cdot$ 2cL with c approx log2 3 characterises the density of the
  Trinity lattice at level L. The factor log2 3 reflects the 3-fold branching of the Lucas
  recurrence --- yet another manifestation of the Trinity anchor identity (star).

  3.3 Concatenation and Composition Rules

        Definition 3.3 (Trinity product). For labels lambda = (n,k,p,m,q) and lambda' =
        (n',k',p',m',q'),

  T(lambda) $\cdot$ T(lambda') = T(nn', k+k', p+p', m+m', q+q').

  This defines a monoid structure on T with identity T(1,0,0,0,0) = 1. The monoid is in
  fact a group if n = 1 (i.e., if the integer prefactor is 1), since then the label lives in Z5
  and the group operation is coordinatewise addition.

        Definition 3.4 (Trinity power). For r $\in$ Q, the fractional power T(lambda)r
        corresponds to the label (nr, rk, rp, rm, rq), valid when nr remains in the allowed
        range and exponents are integers or half-integers.

  3.4 Density and Approximation Properties

        Theorem 3.5 (Informal density statement; conditional). Under the working
        independence hypothesis, for any x > 0 and $\varepsilon$ > 0, there exists T $\in$ T such that |x -
        T| < $\varepsilon$.

  Proof sketch. The group L is a dense subgroup of R under the independence hypothesis
  (since four linearly independent generators over Q generate a dense subgroup of R).
  The exponential of a dense subgroup of R is dense in (0, $\in$fty). □

        Corollary 3.6 (Existence of complexity-minimal approximant; conditional).
        For any real x > 0 and precision $\varepsilon$ > 0, there exists a complexity-minimal Trinity
        monomial T^ with |x - T^| < $\varepsilon$; moreover, the optimal T^* is generically unique
        (uniqueness fails on a set of measure zero if x $\in$ T, and at isolated special values
        otherwise).

  4. The Additive Pellis Expansion

\subsection{References}

- Vasilev, D. \textit{Flos Aureus Monograph: Trinity Constants}. DOI \href{https://zenodo.org/doi/10.5281/zenodo.19227877}{10.5281/zenodo.19227877}.
- CODATA 2018 recommended values: NIST \href{https://physics.nist.gov/cuu/Constants/}{https://physics.nist.gov/cuu/Constants/}.
- Heyrovska, R. Golden-angle FSC publications --- Heyrovský Institute of Physical Chemistry.
